% =====================================================================
% CODEX ColombiaMacro · Parte V · Referencia (español)
% Fragmento para \input: sin preámbulo. Las referencias se agregan al final (05_referencia.tex).
% =====================================================================
\part{Referencia}

% ---------------------------------------------------------------------
\chapter{Limitaciones de los datos}
\label{es:r:limitaciones}

Conocer los límites de cada cifra es parte de leerla bien. Esta es la lista completa de limitaciones conocidas,
agrupadas por tema.

\section{Generales}
\begin{enumerate}
\item \textbf{Una sola versión de la historia.} Cada actualización reemplaza la historia con la última publicación de
  cada entidad. Las cifras pasadas son las revisadas, no las que se conocían en su momento.
\item \textbf{Datos provisionales.} Los últimos trimestres del PIB y las cuentas departamentales son provisionales;
  las cuentas departamentales llegan con más de un año de rezago.
\item \textbf{Fechas distintas en una misma página.} Cada fuente se publica en días diferentes; una página puede
  combinar datos de meses distintos. Cada cifra lleva su fecha.
\item \textbf{Volúmenes encadenados.} Los aportes sectoriales y de la demanda no suman exactamente el total; los
  cálculos ``sin Gobierno'' o ``sin minería'' son aproximaciones.
\end{enumerate}

\section{Ciclo y capacidad}
\begin{enumerate}
\item La brecha del producto y la fase del ciclo dependen del método; el sitio muestra el rango entre métodos por esa
  razón.
\item El tratamiento de la pandemia (interpolación entre 2020-T2 y 2021-T2 para estimar la tendencia) es una decisión
  de modelado.
\item El fechado de expansiones y recesiones usa meses posteriores: el último punto de giro puede moverse o desaparecer
  con datos nuevos.
\end{enumerate}

\section{Precios, tasas y curva}
\begin{enumerate}
\item La tasa neutral es una estimación externa con incertidumbre considerable; la clasificación de la postura depende
  de ella.
\item La inflación implícita contiene primas por riesgo inflacionario y por liquidez y el rezago de indexación de la
  UVR; no es una expectativa pura y las primas pueden tener cualquier signo.
\item Las tasas cero cupón de los TES son salidas de un modelo estimado por el Banco (Nelson-Siegel), no precios
  negociados.
\end{enumerate}

\section{Mercados y empresas}
\begin{enumerate}
\item El COLCAP y los índices equiponderado y de 7 Magníficas son índices de precios sin dividendos.
\item La sección ``bolsa por dentro'' usa fuentes auxiliares no oficiales (composición de un fondo que replica el
  índice y precios de un proveedor comercial); aplicar la canasta vigente hacia atrás introduce sesgo de supervivencia.
\item La lista de las 10.000 empresas cambia cada año y excluye bancos y aseguradoras; las variaciones no comparan las
  mismas empresas y la región es la del domicilio, no la de producción.
\item Las cancelaciones del registro mercantil incluyen depuraciones administrativas; el último mes puede estar
  incompleto.
\item La volatilidad de algunas monedas se calcula con series que repiten el último valor en fines de semana y
  festivos, lo que puede subestimarla frente a series de días hábiles.
\end{enumerate}

\section{Cuentas externas, fiscales y comercio}
\begin{enumerate}
\item Las balanzas comercial y bilateral restan importaciones CIF de exportaciones FOB: el déficit resulta algo mayor
  que en la balanza de pagos.
\item El volumen de comercio es un residuo entre el valor del DANE y los índices de precios del Banco.
\item Las cuentas del Gobierno son de caja: la serie de intereses puede tener meses negativos y subestimar el costo de
  la deuda frente a las cifras en causación.
\item Las razones al PIB en dólares dependen de la tasa de cambio usada para convertir el PIB.
\end{enumerate}

\section{Mercado laboral}
\begin{enumerate}
\item La serie de informalidad empieza en 2021 (nuevo marco muestral) y no es comparable con la anterior.
\item Las series de detalle (ramas, ciudades, jóvenes, sexo) no están desestacionalizadas y tienen más error de
  muestreo que el total nacional.
\end{enumerate}

% ---------------------------------------------------------------------
\chapter{Glosario}
\label{es:r:glosario}

\small
\begin{longtable}{>{\raggedright\arraybackslash}p{3.2cm}>{\raggedright\arraybackslash\itshape}p{3.2cm}>{\raggedright\arraybackslash}p{8.2cm}}
\toprule
\textbf{Término} & \textbf{\upshape En inglés} & \textbf{Definición}\\
\midrule
\endfirsthead
\toprule
\textbf{Término} & \textbf{\upshape En inglés} & \textbf{Definición}\\
\midrule
\endhead
\bottomrule
\endfoot
Amplitud (difusión) & Breadth (diffusion) & Proporción de componentes (sectores, subclases) que cumplen una condición, por ejemplo estar sobre su tendencia.\\
Año corrido & Year to date & Período desde enero hasta el último mes disponible, comparado con el mismo período del año anterior.\\
Apertura comercial & Trade openness & (Exportaciones + importaciones de bienes) / PIB.\\
Apreciación & Appreciation & Aumento del valor de una moneda: se necesitan menos pesos por dólar.\\
Aporte al crecimiento & Contribution to growth & Parte del crecimiento total explicada por un componente: su crecimiento por su peso.\\
Asalariado & Wage employee & Ocupado que trabaja para un empleador a cambio de un salario.\\
Balance primario & Primary balance & Balance fiscal sin el pago de intereses de la deuda.\\
Balanza cambiaria & FX balance & Registro de las divisas que entran y salen por el mercado cambiario regulado.\\
Balanza comercial & Trade balance & Exportaciones menos importaciones de bienes.\\
Balanza de pagos & Balance of payments & Registro de todas las transacciones entre residentes y no residentes en un período.\\
Base 100 & Index base 100 & Índice que divide cada valor por el de una fecha de referencia y multiplica por 100.\\
Bienes de capital & Capital goods & Maquinaria y equipo usados para producir otros bienes; anticipan inversión.\\
Brecha del producto & Output gap & Diferencia porcentual entre el PIB observado y su capacidad (potencial).\\
Breakeven (inflación implícita) & Breakeven inflation & Inflación que iguala el rendimiento de un TES en pesos y uno en UVR del mismo plazo.\\
Capacidad (PIB potencial) & Potential output & Producción sostenible sin presiones sobre los precios; se estima como tendencia.\\
Cartera & Loan book & Saldo de los créditos otorgados por el sistema financiero.\\
CDT & Certificate of deposit & Certificado de depósito a término: ahorro a plazo fijo en un banco.\\
CIF & CIF & Valor de una importación con costo, seguro y flete hasta el puerto de destino.\\
Ciclo económico & Business cycle & Fluctuaciones de la actividad alrededor de su tendencia.\\
COLCAP & COLCAP & Índice de precios de las acciones más líquidas de la Bolsa de Valores de Colombia (MSCI COLCAP desde 2021).\\
Cobertura (comercio) & Export/import coverage & Exportaciones / importaciones, en porcentaje.\\
Cuenta corriente & Current account & Balance de bienes, servicios, ingresos primarios y transferencias con el exterior.\\
Cuenta financiera & Financial account & Flujos de inversión y préstamos entre residentes y no residentes.\\
CUODE & CUODE & Clasificación de las importaciones por uso económico: consumo, materias primas y bienes de capital.\\
Curva de rendimientos & Yield curve & Tasas de los bonos del Gobierno para distintos plazos en una misma fecha.\\
Curva invertida & Inverted curve & Curva en la que el plazo largo rinde menos que el corto.\\
Deflactor & Deflator & Índice de precios implícito que convierte valores nominales en reales.\\
Depreciación & Depreciation & Pérdida de valor de una moneda: se necesitan más pesos por dólar.\\
Desestacionalizar & Seasonal adjustment & Quitar a una serie los patrones que se repiten cada año en las mismas fechas.\\
Desempleo (tasa) & Unemployment rate & Desocupados / fuerza de trabajo.\\
Dólar global & Broad dollar index & Índice de la Reserva Federal del valor del dólar frente a una canasta de monedas.\\
DTF & DTF & Tasa promedio de captación de los CDT a 90 días.\\
Efecto base & Base effect & Cambio de una variación anual explicado por lo que ocurrió hace doce meses.\\
Endeudamiento & Leverage & Pasivos / activos.\\
Equiponderado & Equal-weighted & Índice que da el mismo peso a cada acción.\\
Expectativas ancladas & Anchored expectations & Expectativas de inflación de largo plazo cercanas a la meta y poco sensibles a choques.\\
Expansión & Expansion & Fase con brecha positiva y creciente.\\
Contracción & Contraction & Fase con brecha negativa y decreciente.\\
Desaceleración & Slowdown & Fase con brecha positiva pero decreciente.\\
Recuperación & Recovery & Fase con brecha negativa pero creciente.\\
FOB & FOB & Valor de una exportación puesta a bordo en el puerto de salida, sin fletes ni seguros internacionales.\\
Forward 5y5y & 5y5y forward & Inflación implícita promedio entre los años 5 y 10.\\
Fuerza de trabajo & Labour force & Ocupados más desocupados.\\
GEIH & GEIH & Gran Encuesta Integrada de Hogares del DANE: fuente del mercado laboral.\\
Gobierno nacional central (GNC) & Central government & Nación y sus entidades del orden central, sin empresas públicas ni gobiernos locales.\\
Herfindahl (índice) & Herfindahl index & Suma de las participaciones al cuadrado; mide concentración.\\
IBR & IBR & Indicador Bancario de Referencia: tasa interbancaria de corto plazo.\\
IED & FDI & Inversión extranjera directa: inversión con control o influencia duradera en empresas residentes.\\
Inflación & Inflation & Aumento generalizado de los precios; se mide con el IPC.\\
Inflación de fondo (básica) & Core inflation & Inflación sin los rubros más volátiles (alimentos y regulados).\\
Informalidad & Informality & Ocupados sin protección social o en unidades no registradas, según la definición del DANE.\\
IPC & CPI & Índice de precios al consumidor.\\
ISE & ISE & Indicador de Seguimiento a la Economía: actividad mensual del DANE.\\
ITCR & Real effective exchange rate & Índice de tasa de cambio real: precios de los socios en pesos frente a los precios de Colombia.\\
Logaritmo $\times$100 & Log points & Transformación en la que diferencias pequeñas equivalen a variaciones porcentuales.\\
Margen neto & Net margin & Utilidad / ingresos.\\
Meta de inflación & Inflation target & 3\% con un rango de $\pm1$ pp fijado por el Banco de la República.\\
Nominal & Nominal & Medido en pesos corrientes, sin descontar la inflación.\\
Número equivalente & Equivalent number & 1 / Herfindahl: cuántos componentes iguales darían la misma concentración.\\
Ocupados & Employed & Personas que trabajaron al menos una hora remunerada en la semana de referencia.\\
Pendiente de la curva & Curve slope & TES a 10 años menos TES a 1 año.\\
Percentil & Percentile & Proporción de observaciones históricas menores o iguales al dato actual.\\
PIB & GDP & Producto interno bruto: valor de los bienes y servicios finales producidos en el país.\\
PII & IIP & Posición de inversión internacional: activos menos pasivos externos del país.\\
Postura monetaria & Monetary stance & Tasa real frente a la neutral: restrictiva, neutral o expansiva.\\
pp & pp & Puntos porcentuales: diferencia entre dos porcentajes.\\
Prima por plazo & Term premium & Compensación adicional por mantener un bono de largo plazo.\\
Real & Real & Medido descontando la inflación.\\
Rebase & Rebasing & Volver a dividir un índice para que valga 100 en otra fecha.\\
Regulados & Administered prices & Precios fijados o supervisados por el Estado: energía, gas, transporte, combustibles.\\
Remesas & Remittances & Transferencias de trabajadores en el exterior a sus hogares en Colombia.\\
Rentabilidad del patrimonio & Return on equity & Utilidad / patrimonio.\\
Reservas internacionales & International reserves & Activos externos líquidos en poder del Banco de la República.\\
Subutilización laboral & Labour underutilisation & Desocupados, subocupados y fuerza potencial sobre la fuerza de trabajo ampliada.\\
Suma de 12 meses & 12-month sum & Total de los últimos doce meses; elimina la estacionalidad.\\
Tasa de política (TPM) & Policy rate & Tasa con la que el Banco de la República suministra liquidez a los bancos.\\
Tasa neutral & Neutral rate & Tasa real que ni estimula ni frena la economía.\\
Tasa real ex ante & Ex ante real rate & Tasa nominal descontada la inflación esperada.\\
Tasa real ex post & Ex post real rate & Tasa nominal descontada la inflación observada.\\
Términos de intercambio & Terms of trade & Precio de las exportaciones / precio de las importaciones.\\
TES & TES & Títulos de deuda pública interna del Gobierno colombiano.\\
Traspaso & Pass-through & Proporción del cambio de la tasa de política que llega a otras tasas.\\
TRM & Representative market rate & Tasa de cambio representativa del mercado (pesos por dólar), certificada por la Superfinanciera.\\
Tendencia & Trend & Componente de largo plazo de una serie, estimado con un filtro.\\
UVR & UVR & Unidad de valor real: unidad indexada a la inflación; los TES UVR pagan tasa real.\\
Variación anual & Year-on-year change & Cambio frente al mismo período del año anterior.\\
Variación trimestral anualizada & Annualised quarterly change & Cambio de un trimestre llevado a ritmo anual: $(1+g)^4-1$.\\
Volatilidad & Volatility & Desviación estándar de los cambios diarios, anualizada.\\
Volumen implícito & Implied volume & Cambio del valor descontado el cambio de precio.\\
\end{longtable}
\normalsize

% ---------------------------------------------------------------------
\chapter{Siglas}
\label{es:r:siglas}

\begin{center}
\small
\begin{tabularx}{\linewidth}{lX}
\toprule
\textbf{Sigla} & \textbf{Significado}\\
\midrule
BanRep & Banco de la República\\
BVC & Bolsa de Valores de Colombia\\
CARF & Comité Autónomo de la Regla Fiscal\\
CIIU & Clasificación Industrial Internacional Uniforme\\
DANE & Departamento Administrativo Nacional de Estadística\\
DIAN & Dirección de Impuestos y Aduanas Nacionales\\
EIA & Administración de Información Energética de los Estados Unidos\\
FMI & Fondo Monetario Internacional\\
FRED & Base de datos económicos de la Reserva Federal de San Luis\\
GEIH & Gran Encuesta Integrada de Hogares\\
GNC & Gobierno nacional central\\
HP & Filtro de Hodrick y Prescott\\
IBR & Indicador Bancario de Referencia\\
IED & Inversión extranjera directa\\
IPC & Índice de precios al consumidor\\
ISE & Indicador de Seguimiento a la Economía\\
ITCR & Índice de tasa de cambio real\\
MBP6 & Manual de balanza de pagos, sexta edición (FMI)\\
OCDE & Organización para la Cooperación y el Desarrollo Económicos\\
PIB & Producto interno bruto\\
PII & Posición de inversión internacional\\
RUES & Registro Único Empresarial y Social\\
SPNF & Sector público no financiero\\
TES & Títulos de tesorería\\
TIB & Tasa interbancaria\\
TPM & Tasa de política monetaria\\
TRM & Tasa representativa del mercado\\
UVR & Unidad de valor real\\
\bottomrule
\end{tabularx}
\end{center}

% ---------------------------------------------------------------------
\chapter{Calendario de las fuentes}
\label{es:r:calendario}

El calendario exacto lo fija cada entidad cada año. La tabla resume el patrón habitual de publicación, útil para saber
cuándo esperar un dato nuevo en el sitio.

\begin{center}
\small
\begin{tabularx}{\linewidth}{>{\raggedright\arraybackslash}p{4.2cm}>{\raggedright\arraybackslash}p{2.4cm}X}
\toprule
\textbf{Dato} & \textbf{Frecuencia} & \textbf{Publicación habitual}\\
\midrule
PIB (producción y gasto) & Trimestral & DANE, unos 45 días después del cierre del trimestre\\
ISE & Mensual & DANE, unos 50 días después del mes\\
IPC & Mensual & DANE, primeros días hábiles del mes siguiente\\
GEIH (desempleo, informalidad) & Mensual & DANE, a fin del mes siguiente\\
Exportaciones e importaciones & Mensual & DANE, unos 45–60 días después del mes\\
PIB departamental & Anual & DANE, con más de un año de rezago\\
Tasa de política & Por decisión & Banco de la República, reuniones de la Junta Directiva\\
TRM, IBR, curva TES, COLCAP & Diaria & Banco de la República / Superfinanciera, días hábiles\\
Balanza de pagos & Trimestral & Banco de la República, cerca de 75–90 días después del trimestre\\
Deuda externa, reservas, remesas & Mensual & Banco de la República\\
Balance fiscal del Gobierno (caja) & Mensual & Banco de la República con datos del Ministerio de Hacienda\\
10.000 empresas más grandes & Anual & Superintendencia de Sociedades, a mediados del año siguiente\\
Registro mercantil & Mensual & Confecámaras (RUES), portal de datos abiertos\\
\bottomrule
\end{tabularx}
\end{center}

\chapter{Referencias}
\label{es:r:referencias}
\begin{itemize}[leftmargin=1.2em,itemindent=-1.2em]
\small
\item[] Arango, L. E., Melo, L. F. y Vásquez, D. M. (2002). \emph{Estimación de la estructura a plazo de las tasas de interés en Colombia} (Borradores de Economía n.º 196). Banco de la República.
\item[] Banco de la República. (2025). \emph{Informe de Política Monetaria} (enero, abril y octubre de 2025) y estimaciones de la tasa de interés neutral del equipo técnico. Banco de la República.
\item[] Banco de la República. (s. f.). \emph{Definiciones de la tasa de política monetaria, IBR, TIB, DTF, CDT y tasas de colocación} (catálogo de series estadísticas). Banco de la República.
\item[] Banco de la República. (s. f.). \emph{Mecanismos de transmisión de la política monetaria en Colombia}. Banco de la República.
\item[] Banco de la República. (s. f.). \emph{Metodología de cálculo del Índice de Tasa de Cambio Real (ITCR) de Colombia}. Banco de la República.
\item[] Banco de la República. (s. f.). \emph{Reporte de Estabilidad Financiera} (semestral). Banco de la República.
\item[] Banco de la República, Junta Directiva. (2018). \emph{Resolución Externa 1 de 2018} y Circular Reglamentaria DOAM-146. Banco de la República.
\item[] Baumol, W. J. (1967). Macroeconomics of unbalanced growth: The anatomy of urban crisis. \emph{American Economic Review, 57}(3), 415–426.
\item[] Baxter, M. y King, R. G. (1999). Measuring business cycles: Approximate band-pass filters for economic time series. \emph{Review of Economics and Statistics, 81}(4), 575–593.
\item[] Bell, D. N. F. y Blanchflower, D. G. (2011). Young people and the Great Recession. \emph{Oxford Review of Economic Policy, 27}(2), 241–267.
\item[] Bernanke, B. S. y Gertler, M. (1995). Inside the black box: The credit channel of monetary policy transmission. \emph{Journal of Economic Perspectives, 9}(4), 27–48.
\item[] Bernanke, B. S., Gertler, M. y Gilchrist, S. (1999). The financial accelerator in a quantitative business cycle framework. En J. B. Taylor y M. Woodford (Eds.), \emph{Handbook of Macroeconomics} (Vol. 1C, pp. 1341–1393). Elsevier.
\item[] Betancourt, R., Vargas, H. y Rodríguez, N. (2008). Interest rate pass-through in Colombia: A micro-banking perspective. \emph{Cuadernos de Economía, 45}(131), 29–58.
\item[] Beveridge, S. y Nelson, C. R. (1981). A new approach to decomposition of economic time series into permanent and transitory components with particular attention to measurement of the "business cycle". \emph{Journal of Monetary Economics, 7}(2), 151–174.
\item[] Blanchard, O. (2019). Public debt and low interest rates. \emph{American Economic Review, 109}(4), 1197–1229.
\item[] Blanchard, O., Cerutti, E. y Summers, L. (2015). \emph{Inflation and activity: Two explorations and their monetary policy implications} (IMF Working Paper WP/15/230). Fondo Monetario Internacional.
\item[] Bry, G. y Boschan, C. (1971). \emph{Cyclical analysis of time series: Selected procedures and computer programs}. National Bureau of Economic Research.
\item[] Bryan, M. F. y Cecchetti, S. G. (1994). Measuring core inflation. En N. G. Mankiw (Ed.), \emph{Monetary Policy} (pp. 195–215). University of Chicago Press.
\item[] Burns, A. F. y Mitchell, W. C. (1946). \emph{Measuring business cycles}. National Bureau of Economic Research.
\item[] Calvo, G. A., Leiderman, L. y Reinhart, C. M. (1993). Capital inflows and real exchange rate appreciation in Latin America: The role of external factors. \emph{IMF Staff Papers, 40}(1), 108–151.
\item[] Cecchetti, S. G. (1997). Measuring short-run inflation for central bankers. \emph{Federal Reserve Bank of St. Louis Review, 79}(3), 143–155.
\item[] Cerra, V. y Saxena, S. C. (2008). Growth dynamics: The myth of economic recovery. \emph{American Economic Review, 98}(1), 439–457.
\item[] Chavarro, X., Cristiano, D., Gómez, J. E., González, E. y Huertas, C. (2015). \emph{Evaluación de la transmisión de la tasa de interés de referencia a las tasas de interés del sistema financiero} (Borradores de Economía n.º 874). Banco de la República.
\item[] Chen, Y.-C. y Rogoff, K. (2003). Commodity currencies. \emph{Journal of International Economics, 60}(1), 133–160.
\item[] Christiano, L. J. y Fitzgerald, T. J. (2003). The band pass filter. \emph{International Economic Review, 44}(2), 435–465.
\item[] Clarida, R., Galí, J. y Gertler, M. (1999). The science of monetary policy: A New Keynesian perspective. \emph{Journal of Economic Literature, 37}(4), 1661–1707.
\item[] Congreso de la República de Colombia. (2011). \emph{Ley 1473 de 2011, por medio de la cual se establece una regla fiscal y se dictan otras disposiciones}.
\item[] Congreso de la República de Colombia. (2013). \emph{Ley Estatutaria 1622 de 2013, por medio de la cual se expide el estatuto de ciudadanía juvenil}.
\item[] Congreso de la República de Colombia. (2021). \emph{Ley 2155 de 2021, por medio de la cual se expide la Ley de Inversión Social} (modifica la regla fiscal y crea el Comité Autónomo de la Regla Fiscal).
\item[] DANE. (2019). \emph{Índice de Precios al Consumidor, base diciembre de 2018: metodología y canasta}. Departamento Administrativo Nacional de Estadística.
\item[] DANE. (2021). \emph{Gran Encuesta Integrada de Hogares: metodología y marco conceptual}. Departamento Administrativo Nacional de Estadística.
\item[] DANE. (s. f.). \emph{Cuentas nacionales departamentales, base 2015: metodología (CD-01)}. Departamento Administrativo Nacional de Estadística.
\item[] DANE. (s. f.). \emph{Cuentas nacionales trimestrales: metodología (enfoque de la producción, índices encadenados)}. Departamento Administrativo Nacional de Estadística.
\item[] DANE. (s. f.). \emph{Estadísticas de comercio internacional: exportaciones (FOB) e importaciones (CIF) con base en registros administrativos de la DIAN}. Departamento Administrativo Nacional de Estadística.
\item[] DANE. (s. f.). \emph{Proyecciones y retroproyecciones de población, Censo Nacional de Población y Vivienda 2018 (actualización 2025)}. Departamento Administrativo Nacional de Estadística.
\item[] Dornbusch, R. (1976). Expectations and exchange rate dynamics. \emph{Journal of Political Economy, 84}(6), 1161–1176.
\item[] Espinosa, J. A., Melo, L. F. y Moreno, J. F. (2014). \emph{Estimación de la prima por vencimiento de los TES en pesos del gobierno colombiano} (Borradores de Economía n.º 854). Banco de la República.
\item[] Estrella, A. y Hardouvelis, G. A. (1991). The term structure as a predictor of real economic activity. \emph{Journal of Finance, 46}(2), 555–576.
\item[] Fisher, I. (1930). \emph{The theory of interest}. Macmillan.
\item[] Fondo Monetario Internacional. (2009). \emph{Manual de balanza de pagos y posición de inversión internacional} (6.ª ed., MBP6). FMI.
\item[] Fondo Monetario Internacional, Organización Internacional del Trabajo, OCDE, Eurostat, Comisión Económica de las Naciones Unidas para Europa y Banco Mundial. (2004). \emph{Consumer price index manual: Theory and practice}. FMI.
\item[] Gabaix, X. (2011). The granular origins of aggregate fluctuations. \emph{Econometrica, 79}(3), 733–772.
\item[] Goldin, C. (2014). A grand gender convergence: Its last chapter. \emph{American Economic Review, 104}(4), 1091–1119.
\item[] Gürkaynak, R. S., Sack, B. y Wright, J. H. (2010). The TIPS yield curve and inflation compensation. \emph{American Economic Journal: Macroeconomics, 2}(1), 70–92.
\item[] Hamilton, J. D. (2018). Why you should never use the Hodrick-Prescott filter. \emph{Review of Economics and Statistics, 100}(5), 831–843.
\item[] Hausmann, R., Hwang, J. y Rodrik, D. (2007). What you export matters. \emph{Journal of Economic Growth, 12}(1), 1–25.
\item[] Herfindahl, O. C. (1950). \emph{Concentration in the steel industry} (tesis doctoral). Columbia University.
\item[] Hidalgo, C. A., Klinger, B., Barabási, A.-L. y Hausmann, R. (2007). The product space conditions the development of nations. \emph{Science, 317}(5837), 482–487.
\item[] Hirschman, A. O. (1964). The paternity of an index. \emph{American Economic Review, 54}(5), 761.
\item[] Hodrick, R. J. y Prescott, E. C. (1997). Postwar U.S. business cycles: An empirical investigation. \emph{Journal of Money, Credit and Banking, 29}(1), 1–16.
\item[] Hsieh, C.-T. y Klenow, P. J. (2009). Misallocation and manufacturing TFP in China and India. \emph{Quarterly Journal of Economics, 124}(4), 1403–1448.
\item[] Jaravel, X. (2021). Inflation inequality: Measurement, causes, and policy implications. \emph{Annual Review of Economics, 13}, 599–629.
\item[] Kohli, U. (2004). Real GDP, real domestic income, and terms-of-trade changes. \emph{Journal of International Economics, 62}(1), 83–106.
\item[] Lane, P. R. y Milesi-Ferretti, G. M. (2007). The external wealth of nations mark II: Revised and extended estimates of foreign assets and liabilities, 1970–2004. \emph{Journal of International Economics, 73}(2), 223–250.
\item[] Laubach, T. y Williams, J. C. (2003). Measuring the natural rate of interest. \emph{Review of Economics and Statistics, 85}(4), 1063–1070.
\item[] Litterman, R. y Scheinkman, J. (1991). Common factors affecting bond returns. \emph{Journal of Fixed Income, 1}(1), 54–61.
\item[] Maloney, W. F. (2004). Informality revisited. \emph{World Development, 32}(7), 1159–1178.
\item[] Meese, R. A. y Rogoff, K. (1983). Empirical exchange rate models of the seventies: Do they fit out of sample? \emph{Journal of International Economics, 14}(1–2), 3–24.
\item[] Melitz, M. J. (2003). The impact of trade on intra-industry reallocations and aggregate industry productivity. \emph{Econometrica, 71}(6), 1695–1725.
\item[] Modigliani, F. y Miller, M. H. (1958). The cost of capital, corporation finance and the theory of investment. \emph{American Economic Review, 48}(3), 261–297.
\item[] MSCI. (2021). \emph{MSCI COLCAP Index methodology}. MSCI Inc.
\item[] Naciones Unidas. (1999). \emph{Classification of individual consumption according to purpose (COICOP)}. Naciones Unidas.
\item[] Naciones Unidas, Comisión Europea, FMI, OCDE y Banco Mundial. (2009). \emph{System of National Accounts 2008}. Naciones Unidas.
\item[] Nelson, C. R. y Siegel, A. F. (1987). Parsimonious modeling of yield curves. \emph{Journal of Business, 60}(4), 473–489.
\item[] Obstfeld, M. y Rogoff, K. (1995). The intertemporal approach to the current account. En G. M. Grossman y K. Rogoff (Eds.), \emph{Handbook of International Economics} (Vol. 3, pp. 1731–1799). Elsevier.
\item[] OECD. (2001). \emph{Measuring productivity: OECD manual}. OECD Publishing.
\item[] OECD. (s. f.). \emph{Business Cycle Clock} (metodología del sistema de indicadores adelantados compuestos). OECD.
\item[] Okun, A. M. (1962). Potential GNP: Its measurement and significance. En \emph{Proceedings of the Business and Economic Statistics Section} (pp. 98–104). American Statistical Association.
\item[] Organización Internacional del Trabajo. (1993). \emph{Resolución sobre la Clasificación Internacional de la Situación en el Empleo (CISE-93)}. 15.ª Conferencia Internacional de Estadísticos del Trabajo.
\item[] Organización Internacional del Trabajo. (2013). \emph{Resolución sobre las estadísticas del trabajo, la ocupación y la subutilización de la fuerza de trabajo}. 19.ª Conferencia Internacional de Estadísticos del Trabajo.
\item[] Orphanides, A. y van Norden, S. (2002). The unreliability of output-gap estimates in real time. \emph{Review of Economics and Statistics, 84}(4), 569–583.
\item[] Prebisch, R. (1950). \emph{The economic development of Latin America and its principal problems}. Naciones Unidas, CEPAL.
\item[] Rajan, R. G. y Zingales, L. (1998). Financial dependence and growth. \emph{American Economic Review, 88}(3), 559–586.
\item[] Ravn, M. O. y Uhlig, H. (2002). On adjusting the Hodrick-Prescott filter for the frequency of observations. \emph{Review of Economics and Statistics, 84}(2), 371–376.
\item[] Reinhart, C. M. y Rogoff, K. S. (2009). \emph{This time is different: Eight centuries of financial folly}. Princeton University Press.
\item[] Rey, H. (2013). Dilemma not trilemma: The global financial cycle and monetary policy independence. En \emph{Proceedings of the Jackson Hole Economic Policy Symposium} (pp. 285–333). Federal Reserve Bank of Kansas City.
\item[] Rogoff, K. (1996). The purchasing power parity puzzle. \emph{Journal of Economic Literature, 34}(2), 647–668.
\item[] Superintendencia de Sociedades. (2026). \emph{Informe de las 10.000 empresas más grandes del país} (cifras a diciembre de 2025). Supersociedades.
\item[] Superintendencia Financiera de Colombia. (s. f.). \emph{Tasa de cambio representativa del mercado: antecedentes normativos y metodología}. Superfinanciera.
\item[] Taylor, J. B. (1993). Discretion versus policy rules in practice. \emph{Carnegie-Rochester Conference Series on Public Policy, 39}, 195–214.
\end{itemize}